% Section: a statistical theory of the daily bar under a noisy open.
% Every number is \thn{key}, generated by scripts/45_theory_checks.py.
\section{A statistical theory of the daily bar under a noisy open}\label{sec:theory}

Sections~6.6--6.8 of the paper report four facts:
\begin{itemize}
  \item NEPSE's open is mostly reversed during the session;
  \item a wider pre-open band lowered the unbiasedness coefficient sharply;
  \item an estimator that ignores the open forecast better than one that reads it, in every
        frontier panel;
  \item a pooled calibration lagged the band reform.
\end{itemize}
This section gives the statistics behind those facts in six propositions. Proofs are in
Appendix~\ref{app:proofs}. The script \texttt{scripts/45\_theory\_checks.py} compared each closed
form, identity and inequality with a simulation or a numerical integral, as plan M19 required. The
plan was registered before the script was written. Table~\ref{tab:checks} lists the checks. The
propositions are post hoc relative to plans M14--M18, because they were derived after those
results were known. Only the test of Section~\ref{sec:pooling} had its prediction fixed in advance.

\subsection{The model}\label{sec:model}

For session $t$, let $C_{-}$ be the printed previous close and $O,H,L,C$ the printed open, high,
low and close. The bar's coordinates are:
\begin{itemize}
  \item the overnight return $o=\ln(O/C_{-})$;
  \item the intraday return $c=\ln(C/O)$;
  \item the close-to-close return $r=o+c$;
  \item the high and low relative to the open, $u=\ln(H/O)$ and $d=\ln(L/O)$;
  \item the range $R=u-d$.
\end{itemize}
Following Appendix~A of the paper, write
\begin{equation}\label{eq:model}
  o = e_o + \eta - \varepsilon_{-}, \qquad c = e_c - \eta + \varepsilon, \qquad r = e_o + e_c + \varepsilon - \varepsilon_{-}.
\end{equation}
Here $e_o$ and $e_c$ are the efficient overnight and intraday returns, $\eta$ is the open's
transient error, and $\varepsilon$ and $\varepsilon_{-}$ are the errors of today's and the
previous close. Two combinations do all the work:
\begin{itemize}
  \item $\os=e_o-\varepsilon_{-}$ is the efficient open measured from the printed previous close;
  \item $x=e_c+\varepsilon$ is everything that happens after the open.
\end{itemize}
Then $o=\os+\eta$ and $r=\os+x$. All moments are raw, as in
\begin{equation}\label{eq:b}
  b=\frac{\E[or]}{\E[o^2]}.
\end{equation}
This is the unbiasedness coefficient of the price-discovery literature
\citep{BiaisHillionSpatt1999,BarclayHendershott2003}.

\begin{assumption}[Nothing known at the open predicts the rest of the session]\label{as:orth}
Second moments are finite, $\E[o^2]>0$, and $\E[\os x]=\E[\eta x]=0$.
\end{assumption}

Assumption~\ref{as:orth} holds if efficient prices are martingales with respect to information
that includes the open's error, and if the closing error is not predictable from the open. It
places no restriction on how the error relates to overnight news. The open may overreact
($\E[\os\eta]>0$), underreact ($\E[\os\eta]<0$) or be stale.

\begin{assumption}[A Brownian session with an instantly corrected open]\label{as:path}
\leavevmode
\begin{itemize}
  \item The closing errors are zero.
  \item After the printed open $o=\os+\eta$, the price follows the efficient path
        $\os+\sigma_c W(s)$, $s\in(0,1]$, where $W$ is a standard Brownian motion independent of
        $(\os,\eta)$.
  \item The high and the low are the extremes of the printed path,
        $\ln(H/C_{-})=\max(o,\os+\sigma_c M)$ and $\ln(L/C_{-})=\min(o,\os+\sigma_c m)$, where
        $M=\sup W$ and $m=\inf W$.
  \item The close is efficient, so $x=\sigma_c D$ with $D=W(1)$.
\end{itemize}
Where the error's law matters, $\eta\sim N(0,\sigma_\eta^2)$ independently of $(\os,W)$. The
error's scale relative to the session is $\rho=\sigma_\eta/\sigma_c$.
\end{assumption}

\begin{assumption}[A band on the open]\label{as:band}
The auction would print a latent open $o_L$, with $\E[r\mid o_L]=b_L\,o_L$. The printed open is
$o=g(o_L)=\max\{-\beta,\min(\beta,o_L)\}$, and the band leaves the close unaffected.
\end{assumption}

A linear conditional mean holds for every elliptical law of $(o_L,r)$, Gaussian or Student-$t$
\citep{CambanisHuangSimons1981}. Assumption~\ref{as:band} therefore allows heavy tails.

\subsection{What the unbiasedness coefficient identifies}\label{sec:identification}

\begin{proposition}[Identification]\label{prop:ident}
Under Assumption~\ref{as:orth}:
\begin{enumerate}[label=(\alph*)]
  \item $\E[or]=\E[o\os]$, so $b=\arg\min_\gamma \E[(\os-\gamma o)^2]$. The effective open
        $C_{-}\exp(b\,o)$ is the best linear predictor of the efficient open. If $\E[\os\mid o]$
        is linear in $o$, then $b\,o=\E[\os\mid o]$.
  \item $\E[oc]=-(1-b)\,\E[o^2]$: the session reverses the fraction $1-b$ of the overnight move.
  \item $\E[\eta^2]=(1-b)^2\,\E[o^2]+\E[(\os-b\,o)^2]$. Hence
        \begin{equation}\label{eq:sharp}
          \E[\eta^2]\;\ge\;(1-b)^2\,\E[o^2],
        \end{equation}
        with equality if and only if $\os=b\,o$ almost surely.
  \item Three special cases:
        \begin{itemize}
          \item If $\E[\os\eta]=0$, then $b=\E[\os^2]/(\E[\os^2]+\E[\eta^2])$ and
                $\E[\eta^2]=(1-b)\E[o^2]$.
          \item If $\eta=\theta\os+\nu$ with $\E[\os\nu]=0$, then
                $b=(1+\theta)\E[\os^2]/\{(1+\theta)^2\E[\os^2]+\E[\nu^2]\}$.
          \item If the open moves only part of the way, $o=\lambda\os$, then $b=1/\lambda$.
        \end{itemize}
  \item The second moments of $(o,r)$ identify $b$, $\E[o^2]$ and $\E[r^2]=\E[\os^2]+\E[x^2]$.
        They do not identify how $\E[r^2]$ divides between the open and the session. Every value
        of $\E[\os^2]$ in $[b^2\E[o^2],\,\E[r^2]]$ is attained by some law satisfying
        Assumption~\ref{as:orth} with the same Gaussian law of $(o,r)$. Correspondingly,
        $\E[\eta^2]$ ranges over $[(1-b)^2\E[o^2],\,(1-b)^2\E[o^2]+\E[r^2]-b^2\E[o^2]]$.
\end{enumerate}
\end{proposition}

Part~(c) is a variance decomposition. The error's second moment is the reversal the data show,
$(1-b)^2\E[o^2]$, plus the part of the efficient open that the printed open cannot predict.
The decomposition yields the sharp form of the bound given in Appendix~A of the paper. That
bound, $\E[\eta^2]\ge \E[o^2]\{(\sqrt{5-4b}-1)/2\}^2$, is valid, but it is strictly smaller
than~\eqref{eq:sharp} for every $b<1$. At $b=0.3$ the sharp bound is \thn{p1-ratio-at-03} times
larger. Appendix~A also states that its bound is attained when the error is perfectly correlated
with news. That is incorrect. Under perfect correlation, $\eta=\{(1-b)/b\}\os$, and the error's
second moment equals the sharp bound, $(1-b)^2\E[o^2]$. The correction is registered as M-026.

Part~(e) is why the paper sets the level of every estimator on the close-to-close scale. The only
level the bar identifies without assumptions about the error is $\E[r^2]$. How it divides between
the overnight move and the session is identified only up to the interval in~(e). The overnight
term of Anam's estimator, $(b\,o)^2$, has expectation $b^2\E[o^2]$, the lower end of that
interval. That lower end is exact under pure proportional overreaction. Part~(d) gives another
reading of the S\&P~500's pooled $b$ of 2.21 (Section~6.8): it is what an index open
reflecting about $45\%$ of the overnight move would produce.

\paragraph{Application (plan M19, part B1).}
Table~\ref{tab:bounds} recomputes the bound with the joint bootstrap of plan M15, over the rows
of that plan's post hoc table. Under the $\pm2\%$ band in the first regime, the sharp bound puts
the open's error at no less than \thn{b1-A1-sharp-pct}\% of the open-to-close proxy, with 95\%
interval $[\thn{b1-A1-sharp-lo}, \thn{b1-A1-sharp-hi}]$. This is \thn{b1-A1-sharp-over-pub}
times the published bound of \thn{b1-A1-pub-pct}\%. Under independence the share is
\thn{b1-A1-ind-pct}\%. After the band was widened (regime~C), the sharp bound is
\thn{b1-C-sharp-pct}\%. On the NIFTY~50, $b=\thn{b1-NIFTY-b}$, and the bound is
\thn{b1-NIFTY-sharp} with interval $[\thn{b1-NIFTY-sharp-lo}, \thn{b1-NIFTY-sharp-hi}]$, which is
indistinguishable from zero.

\begin{table}[t]
\centering\small
\caption{The open's error as a share of the open-to-close proxy $\E[c^2]$ (Proposition~\ref{prop:ident}(c)), with 95\% joint-bootstrap intervals (plan M15's bootstrap, 499 replicates).}\label{tab:bounds}
\begin{tabular}{lccccc}
\toprule
Sample & $b$ & $\E[o^2]/\E[c^2]$ & Published bound & Sharp bound & Independence \\
 & & & $\{\tfrac{\sqrt{5-4b}-1}{2}\}^2\tfrac{\E[o^2]}{\E[c^2]}$ & $(1-b)^2\tfrac{\E[o^2]}{\E[c^2]}$ & $(1-b)\tfrac{\E[o^2]}{\E[c^2]}$ \\
\midrule
\tabinput{generated/tab_bounds.tex}
\bottomrule
\end{tabular}
\end{table}

\subsection{What the classical estimators measure}\label{sec:classical}

Write the classical estimators as follows:
\begin{itemize}
  \item Parkinson's, $P=R^2/(4\ln2)$ \citep{Parkinson1980};
  \item Garman and Klass's, $GK=\tfrac12R^2-(2\ln2-1)c^2$ \citep{GarmanKlass1980};
  \item Rogers and Satchell's, $RS=u(u-c)+d(d-c)$ \citep{RogersSatchell1991};
  \item the daily form of Yang and Zhang's, $YZ=o^2+k\,c^2+(1-k)RS$, with $k=\thn{p2-yz-k}$ for a
        21-session window \citep{YangZhang2000}.
\end{itemize}
Let $T_1(\rho)=\{(1+\rho^2)\arctan\rho-\rho\}/\pi$.

\begin{proposition}[Biases under a noisy open]\label{prop:classical}
\leavevmode
\begin{enumerate}[label=(\alph*)]
  \item Under Assumption~\ref{as:orth} alone, $o^2+c^2-r^2=-2oc$. Hence
        $\E[o^2+c^2]-\E[r^2]=2(1-b)\E[o^2]$ whatever the joint law of $(\os,\eta,x)$.
\end{enumerate}
Under Assumption~\ref{as:path}, with $\eta$ independent of $(\os,W)$ and Gaussian:
\begin{enumerate}[label=(\alph*),resume]
  \item Close-to-close is unbiased: $\E[r^2]=\E[\os^2]+\sigma_c^2$. The overnight square and the
        open-to-close square each carry the whole error:
        $\E[o^2]-\E[\os^2]=\E[c^2]-\sigma_c^2=\sigma_\eta^2$.
  \item Parkinson. $\E[R^2]=\sigma_c^2\{4\ln2+\Delta_R(\rho)\}$, where
        \begin{equation}\label{eq:deltaR}
          \Delta_R(\rho)=2T_1(\rho)+\tfrac{4}{\pi}(\rho-\arctan\rho)+4\int_0^\infty J(s)\,\Phibar(s/\rho)\,ds,
        \end{equation}
        with $J(s)=\E[(-m)\ind\{M<s\}]$, an explicit series given in Appendix~\ref{app:proofs}.
        Moreover, $0\le\Delta_R(\rho)\le (8/\pi)\rho+\rho^2$. As $\rho\to0$,
        $\Delta_R(\rho)=2\ln2\,\rho^2\{1+o(1)\}$.
  \item Rogers--Satchell. $\E[RS]-\sigma_c^2=\sigma_c^2\{\rho^2/2-T_1(\rho)\}$, exactly.
  \item Garman--Klass. $\E[GK]-\sigma_c^2=\sigma_c^2\{\Delta_R(\rho)/2-(2\ln2-1)\rho^2\}$.
  \item The probability that the open is the session's high or low is $(2/\pi)\arctan\rho$.
  \item As $\sigma_\eta\to0$ with $\sigma_c$ fixed, each estimator's bias is
        $\gamma\sigma_\eta^2\{1+o(1)\}$. The loading $\gamma$ is:
        \begin{itemize}
          \item 0 for close-to-close;
          \item 1 for $c^2$;
          \item $\tfrac12$ for Parkinson and for Rogers--Satchell;
          \item $1-\ln2=\thn{p2-gk-load}$ for Garman--Klass;
          \item $\tfrac32$ for $o^2+P$;
          \item $2-\ln2=\thn{p2-gk2-load}$ for $o^2+GK$;
          \item $(3+k)/2=\thn{p2-yz-load}$ for the daily Yang--Zhang form.
        \end{itemize}
\end{enumerate}
\end{proposition}

Part~(a) is the Yang--Zhang excess of Section~6.7 of the paper written in terms of $b$. The
overnight and open-to-close squares together double-count by exactly twice the reversal. Part~(c)
is not a first-order approximation in disguise:
\begin{itemize}
  \item the exact $\Delta_R(\rho)/(2\ln2\,\rho^2)$ is \thn{p2-Pratio-0p5} at $\rho=0.5$,
        \thn{p2-Pratio-1p0} at $\rho=1$ and \thn{p2-Pratio-2p0} at $\rho=2$;
  \item Parkinson therefore absorbs about half the error's variance until the error is as large as
        the session's own volatility.
\end{itemize}
Part~(f) quantifies the channel by which Parkinson, which never reads the open directly, inherits
the open's error. The loadings in~(g) rank the estimators by their exposure to the open. Only
close-to-close has none. Garman--Klass offsets part of the range's excess with the open-to-close
square. The estimators that add the whole overnight move carry one and a half to
\thn{p2-yz-load} times the error's variance.

\paragraph{Application (plan M19, part B2).}
Proposition~\ref{prop:classical}(f) can be inverted. Under Assumption~\ref{as:path}, the share $s$
of non-stale opens that are the session's high or low implies $\sigma_\eta/\sigma_c=\tan(\pi s/2)$.
Under independence, $b$ implies $\sqrt{x/(1-x)}$, where $x=(1-b)\E[o^2]/\E[c^2]$ by
Proposition~\ref{prop:ident}(d). In the first regime, \thn{b2-A1-share}\% of non-stale opens are
an extreme. The two estimates of $\sigma_\eta/\sigma_c$ are:
\begin{itemize}
  \item from that share, \thn{b2-A1-arc} $[\thn{b2-A1-arc-lo}, \thn{b2-A1-arc-hi}]$;
  \item from $b$, \thn{b2-A1-ind} $[\thn{b2-A1-ind-lo}, \thn{b2-A1-ind-hi}]$.
\end{itemize}
After the reform the share is \thn{b2-C-share}\%, and the two estimates are \thn{b2-C-arc} and
\thn{b2-C-ind}. The share-based scale is the larger in every regime, as M19 said it would be if
the open is an extreme more often than an instantly corrected Gaussian error explains. A thin
book makes the first trade the session's extreme even without an error, and a slowly corrected
error keeps the price near the open for longer. Both readings agree on the direction of the
reform: the error grew relative to the session.

\subsection{What a price band does to the unbiasedness coefficient}\label{sec:band}

\begin{proposition}[Censoring by a band]\label{prop:band}
Under Assumption~\ref{as:band}:
\begin{enumerate}[label=(\alph*)]
  \item $b=b_L\,F$, with $F=\E[g(o_L)o_L]/\E[g(o_L)^2]\ge1$. Equality holds if and only if
        $P(|o_L|>\beta)=0$.
  \item On interior opens ($|o_L|<\beta$), the coefficient equals $b_L$. On opens pinned at the
        band, it equals $b_L\,\E[|o_L|\mid|o_L|\ge\beta]/\beta\ge b_L$.
  \item If $o_L\sim N(0,\sigma^2)$ and $k=\beta/\sigma$, then
        \begin{equation}\label{eq:F}
          F(k)=\frac{2\Phi(k)-1}{2\Phi(k)-1-2k\phi(k)+2k^2\Phibar(k)}.
        \end{equation}
        $F$ is strictly decreasing, with $kF(k)\to2\phi(0)$ as $k\to0$ and $F(k)\to1$ as
        $k\to\infty$. The pinned-to-interior ratio is $\lambda(k)/k$, where
        $\lambda=\phi/\Phibar$ is the inverse Mills ratio.
  \item For two regimes $0$ and $1$,
        $\ln(b_0/b_1)=\ln(b_0/b_{0}^{\mathrm{int}})+\ln(b_{0}^{\mathrm{int}}/b_{1}^{\mathrm{int}})+\ln(b_{1}^{\mathrm{int}}/b_1)$.
        Under Assumption~\ref{as:band} in both regimes:
        \begin{itemize}
          \item the first term is $\ln F_0>0$ (censoring);
          \item the second is the change in the latent coefficient;
          \item the third is $-\ln F_1\le0$.
        \end{itemize}
        A positive third term contradicts a linear $\E[r\mid o_L]$ in regime~1.
\end{enumerate}
\end{proposition}

The numerator in~\eqref{eq:F} is Stein's identity: $\E[g(z)z]=\E[g'(z)]$ is the probability of
an interior open. A tight band raises the measured coefficient mechanically, because a pinned open
understates the move the close completes.

\paragraph{Application (plan M19, part B3).}
Table~\ref{tab:censoring} sets the Gaussian factor, computed from each regime's pinned share,
against the ratio the data show. Under the $\pm2\%$ band:
\begin{itemize}
  \item between \thn{b3-B-pin} and \thn{b3-A1-pin} of non-stale opens were pinned, which puts the
        band \thn{b3-A1-k}--\thn{b3-B-k} standard deviations out;
  \item the factor is \thn{b3-B-F}--\thn{b3-A1-F};
  \item the observed $b/b^{\mathrm{int}}$ is \thn{b3-A1-Fobs} in A1, \thn{b3-B-Fobs} in B and
        \thn{b3-A2-Fobs} in A2.
\end{itemize}
The observed ratios fall on both sides of the Gaussian values. The censoring factor and the
pinned-to-interior ratio are below them in A1 and B, and above them in A2. The Gaussian law
therefore gives the order of magnitude of the band's effect, not its exact size.

Part~(d) decomposes the fall of $b$ at the band reform, from \thn{b3-A2-b} in A2 to
\thn{b3-C-b} in C (a factor of \thn{b3-dec-ratio}), into three terms:
\begin{itemize}
  \item the end of censoring, \thn{b3-dec-cens} or \thn{b3-dec-cens-share}\% of the log fall;
  \item a change in the interior coefficient, \thn{b3-dec-int} or \thn{b3-dec-int-share}\%: the
        interior coefficient was \thn{b3-A2-bint} before and \thn{b3-C-bint} after;
  \item a third term of \thn{b3-dec-nonlin}, or \thn{b3-dec-nonlin-share}\%.
\end{itemize}
The third term is positive. By~(d) that is impossible if $\E[r\mid o_L]$ is linear in C. Opens
that the old band would have pinned were reversed more than interior ones after the widening
($b=\thn{b3-C-b}$ overall against \thn{b3-C-bint} inside $\pm1.9\%$). This is the overshoot
recorded in Section~6.7. On this decomposition, about a third of the fall in $b$ is the band's
mechanical inflation coming to an end. Almost two thirds is the overshoot of the large opens the
band used to cap.

\begin{table}[t]
\centering\small
\caption{The censoring factor of Proposition~\ref{prop:band} against NEPSE's regimes. The pinned share is among non-stale opens. The ratios are arithmetic on plan M15's Table~89.}\label{tab:censoring}
\begin{tabular}{lcccccc}
\toprule
Regime & Pinned share & $k$ & $F(k)$ & Observed $b/b^{\mathrm{int}}$ & $\lambda(k)/k$ & Observed $b^{\mathrm{pin}}/b^{\mathrm{int}}$ \\
\midrule
\tabinput{generated/tab_censoring.tex}
\bottomrule
\end{tabular}
\end{table}

\subsection{One coefficient pooled across securities whose opens differ}\label{sec:pooling}

Securities $i$ are forecast with $f_{it}$ for a target $Y_{it}$, the mean squared close-to-close
return over the next $h$ sessions. The QLIKE loss is
$\QL(Y,f)=Y/f-\ln(Y/f)-1$, which ranks forecasts as the latent variance would
\citep{Patton2011}. Write $c^*_i=\E_i[Y/f]$, the mean over security $i$'s forecasts, and
$q(c)=c-1-\ln c\ge0$.

\begin{proposition}[Pooling]\label{prop:pool}
\leavevmode
\begin{enumerate}[label=(\alph*)]
  \item The shape-plus-level identity, for every security:
        $\E_i[\QL(Y,f)]=\E_i[\QL(Y,c^*_if)]+q(c^*_i)$. The first term is the security's
        \emph{shape} loss, the loss left after the best rescaling of its forecasts. The second is
        its \emph{level} loss. The identity holds exactly for sample means.
  \item A common rescaling $\kappa$ of all forecasts that minimises the pooled loss sets the
        row-weighted mean of the $c^*_i$ to one. The pooled level loss is then
        $\sum_iw_iq(c^*_i)\approx\tfrac12\sum_iw_i(\ln c^*_i)^2$, the cross-sectional dispersion
        of the securities' calibration.
  \item Suppose Assumptions~\ref{as:orth}--\ref{as:path} hold for each security, with independent
        errors. For the kernel $K(\beta)=(\beta o)^2+R^2/(4\ln2)$, the level ratio
        $\psi_i(\beta)=\E[K_i(\beta)]/\E[r_i^2]$ is
        $1-(1-\beta^2)s_i+\beta^2n_i+\sigma_{c,i}^2\Delta_R(\rho_i)/(4\ln2\,\E[r_i^2])$, where
        $s_i=\E[\os_i{}^2]/\E[r_i^2]$ and $n_i=\sigma_{\eta,i}^2/\E[r_i^2]$. To first order in the
        error, $\psi_i(\beta)=1-(1-\beta^2)s_i+(\beta^2+\tfrac12)n_i$.
  \item In that first-order form, two cases follow:
        \begin{itemize}
          \item If securities differ only in the error ($s_i\equiv s<2/3$), the cross-sectional
                variance of $\ln\psi_i(\beta)$ is strictly increasing in $\beta\in[0,1]$.
          \item If they differ only in the overnight share, it is strictly decreasing and is zero
                at $\beta=1$.
        \end{itemize}
\end{enumerate}
\end{proposition}

Proposition~\ref{prop:pool} explains when reading the open costs more than it gives.
\begin{itemize}
  \item \textbf{The cost is in the level.} A pooled calibration cannot fit every security, and
        reading the open makes each security's level depend on its own error variance. The level
        loss in~(b) therefore grows with $\beta$ when opens differ.
  \item \textbf{The gain is in the shape.} An informative open tracks the overnight variance.
\end{itemize}
In \thn{p4-sim-nsec} simulated securities sharing one error variance, the full form had the lower
loss (difference \thn{p4-sim-hom-dtotal}). When the error variances were spread over two orders of
magnitude, the open-free form had the lower loss (\thn{p4-sim-het-dtotal}). All of that advantage
was a level advantage: the level difference was \thn{p4-sim-het-dlevel} and the shape difference
\thn{p4-sim-het-dshape}. The security-level loss difference fell with the security's own $b$
(Spearman \thn{p4-sim-het-rho}, against \thn{p4-sim-hom-rho} with a common error).

\paragraph{The test (plan M19, part C).}
Plan M19 fixed a prediction before any of these statistics was computed on data: the open-free
form gains most where a security's own open is noisiest. It also fixed the verdicts. The test uses
the frozen forecasts of the five panel test spans at both horizons, \thn{c-n-cases} cases in all
(Table~\ref{tab:pooling}). In each case the per-security loss difference $d_i$ (full minus
open-free) is correlated with the security's own $b_i$ over the training span.

\begin{itemize}
  \item \textbf{The level component.} The open-free form had the lower loss in \thn{c-n-pos} of
        the \thn{c-n-cases} cases, and its level loss was lower in \thn{c-n-dlevel-pos} of them. The
        level component of part~(b) favours the open-free form in every panel.
  \item \textbf{The mechanism, P4a.} The frozen reading required more: that the full form keep the
        better shape. That held in only \thn{c-n-level} cases. In \thn{c-n-dshape-pos} cases the
        full form's shape loss was higher too, so the mechanism verdict is \emph{\thn{c-p4a}}. The
        open-free form's advantage is therefore not only a matter of calibration: reading a noisy
        open also worsens how well the forecast tracks.
  \item \textbf{The cross-section, P4b.} The prediction was confirmed in \thn{c-n-confirmed}
        cases: NEPSE and Dhaka 2023--2026 at both horizons. It was not detected in
        \thn{c-n-notdetected}: Vietnam and Morocco at both horizons, and Dhaka 2009--2021 at
        twenty-one sessions. It was reversed in \thn{c-n-reversed}: Dhaka 2009--2021 at five
        sessions. The frozen overall verdict is therefore \emph{\thn{c-p4b}}. The ten cases are
        reported without a multiplicity adjustment.
\end{itemize}
The prediction held in the two panels with the lowest pooled coefficient on their training spans,
\thn{c-nep-bpool} and \thn{c-dsea-bpool}. It failed where the coefficient is between
\thn{c-dseb-bpool} and \thn{c-ma-bpool}. That pattern was not predicted. It is reported here and
has not been tested.

\begin{table}[t]
\centering\footnotesize\setlength{\tabcolsep}{4pt}
\caption{The pooling test of plan M19 (part C). $\Delta$: full form minus open-free form, row-weighted mean QLIKE, split by Proposition~\ref{prop:pool}(a). $\rho_S$: Spearman correlation of the security-level loss difference with the security's own $b$ (training span), with a 95\% bootstrap interval over securities (1{,}999 resamples). Predicted: $\rho_S<0$.}\label{tab:pooling}
\begin{tabular}{lrrrrrll}
\toprule
Panel & $h$ & Securities & $\Delta$ total & $\Delta$ level & $\Delta$ shape & $\rho_S$ [95\% interval] & Verdict \\
\midrule
\tabinput{generated/tab_pooling.tex}
\bottomrule
\end{tabular}
\end{table}

\subsection{A rolling calibration after a rule change}\label{sec:lag}

Let $R_s=\sum_ir_{is}^2$ and $A_s=\sum_iA_{is}$ be date-$s$ sums of squared returns and of an
estimator's kernel, and let $\kappa_t=\sum_{s\in W_t}R_s\big/\sum_{s\in W_t}A_s$ over the last $L$
dates. Suppose that before a rule change at $\tau$ the date sums satisfy
$R_s/N-\rho_0a_s\to0$ and $A_s/N-a_s\to0$ in probability as the number of securities $N$ grows,
for a kernel mass $a_s>0$ per security. After the change the same holds with $\rho_1$.

\begin{proposition}[Calibration lag]\label{prop:lag}
\leavevmode
\begin{enumerate}[label=(\alph*)]
  \item As $N\to\infty$, $\kappa_t-\{\omega_t\rho_1+(1-\omega_t)\rho_0\}\to0$ in probability,
        where $\omega_t=\sum_{s\in W_t,\,s\ge\tau}a_s/\sum_{s\in W_t}a_s$ is the new regime's
        share of the window's kernel mass. The relative bias $(1-\omega_t)(\rho_0/\rho_1-1)$
        vanishes only once the window has passed the change.
  \item If the kernel mass is constant across the change, $\omega_{\tau+j}=\min\{(j+1)/L,1\}$:
        the bias decays linearly over $L$ dates.
  \item Suppose changes of log size $\pm\delta$ arrive at rate $\lambda$ per date, and within a
        regime $\ln\kappa_t$ has variance $v/L$. Then, to second order in $\delta$ and for small
        $\lambda L$, the average level loss $\tfrac12\E[(\ln\kappa_t-\ln\rho_t)^2]$ is
        $\tfrac12(\lambda\delta^2L/3+v/L)$. It is minimised at
        $L^*=\sqrt{3v/(\lambda\delta^2)}$, where it equals $\sqrt{\lambda\delta^2v/3}$.
\end{enumerate}
\end{proposition}

Part~(c) is the bias--variance choice of an estimation window under breaks
\citep{PesaranTimmermann2007}, applied to a calibration ratio. The factor $1/3$ is
$\int_0^1(1-u)^2du$, from the linear decay in~(b).

\paragraph{Application (plan M19, part B4).}
The frozen calibration of Anam's estimator on NEPSE is the ratio over the 60 dates before the
band reform, $\rho_0=\thn{b4-rho0}$. Regime~C's own ratio is $\rho_1=\thn{b4-rho1}$. Over C's
\thn{b4-nC} sessions, the path of Proposition~\ref{prop:lag}(a) explains \thn{b4-r2-pct}\% of the
squared deviation of the applied calibration from $\rho_1$ (root mean squared error
\thn{b4-rmse}). The linear weights of part~(b) explain a share of \thn{b4-r2-lin}.

The before-reform regimes give $v=\thn{b4-v}$. NEPSE's three rule boundaries changed the
calibration by $\delta=\thn{b4-delta-A1B}$, \thn{b4-delta-BA2} and \thn{b4-delta-A2C} over
\thn{b4-T} sessions. With these, part~(c) puts the balancing window at $L^*\approx\thn{b4-Lstar}$
dates, shorter than the frozen 60.

Regime~C lies wholly after a change. There, the lag alone, with a constant post-reform ratio,
predicts that a 60-date calibration loses \thn{b4-pred-60} more than a 10-date one. The observed
QLIKE differences are \thn{b4-obs-60-5} at five sessions and \thn{b4-obs-60-21} at twenty-one, so
the lag accounts for about \thn{b4-pred-share-5}\% of the gap. Proposition~\ref{prop:lag}(a)
describes the applied calibration's path almost exactly. The forecast loss of a long window in C,
however, comes mostly from something the proposition holds fixed. The ratio itself moved within C,
because the estimator's own pooled $b$ was still falling and so the kernel was changing, and a
short window tracks that movement. The balancing window $L^*$ assumes the same constant ratios, and
is likely too long for the same reason.

\subsection{Identifying an estimator's calibration without a benchmark}\label{sec:iv}

Plan M14 writes each daily estimator as $X_{k,it}=\alpha_{k,i}+\beta_kIV_{it}+U_{k,it}$, where
$\E[U_{it}\mid IV_{it},\mathcal F_{i,t-1}]=0$. The reference is $X_0=c^2$, with $\beta_0=1$.
Let $V_{it}=\E[IV_{it}\mid\mathcal F_{i,t-1}]$.

\begin{proposition}[Instrumented calibration]\label{prop:iv}
Under this model, for any $\mathcal F_{i,t-1}$-measurable instruments $Z$ with finite second
moments:
\begin{enumerate}[label=(\alph*)]
  \item $\Cov(X_k,Z)=\beta_k\Cov(X_0,Z)$ for every $k$. The cross-covariance matrix has rank one,
        and $\beta_k=\Cov(X_k,Z'w)/\Cov(X_0,Z'w)$ whenever $\Cov(V,Z'w)\neq0$.
  \item With $IV$ uncorrelated with $U$, the least-squares slope of $X_k$ on $X_0$ is
        $\{\beta_k\Var(IV)+\Cov(U_k,U_0)\}/\{\Var(IV)+\Var(U_0)\}$. The mean ratio is
        $\beta_k+\delta_k$, with $\delta_k=(\E X_k-\beta_k\E X_0)/\E X_0$. Neither is $\beta_k$ in
        general.
  \item If an instrument dated $t-\ell$ is correlated with the measurement errors, the
        instrumented estimand is
        $\{\beta_k\Cov(V,Z)+\Cov(U_k,Z)\}/\{\Cov(V,Z)+\Cov(U_0,Z)\}$. If the errors are moving
        averages of order $q$ with respect to $\mathcal F$, instruments dated $t-q-1$ or earlier
        are valid.
  \item The two-stage least-squares estimator with forward orthogonal deviations is consistent and
        asymptotically normal under standard mixing and moment conditions. Hansen's statistic
        is asymptotically $\chi^2$ with one degree of freedom fewer than the number of excluded
        instruments \citep{Hansen1982,ArellanoBover1995}. With moving-average errors, its
        covariance must allow for serial correlation of the same order.
  \item Among composites $w'X$ with $w'\beta=1$, the variance is minimised by
        $w^*=\Sigma_X^{-1}\beta/(\beta'\Sigma_X^{-1}\beta)$. This equals
        $\Omega_U^{-1}\beta/(\beta'\Omega_U^{-1}\beta)$ when $\Sigma_X=\Var(IV)\beta\beta'+\Omega_U$.
\end{enumerate}
\end{proposition}

Part~(c) qualifies the identification behind Section~6.6.
\begin{itemize}
  \item \textbf{Valid instruments.} In \thn{p6-reps} simulated panels with serially independent
        errors, instruments dated $t-1$ recovered slopes of 0.8 and 1.3 as \thn{p6-beta1} and
        \thn{p6-beta2}. The $J$ test rejected at 5\% in a share of \thn{p6-size} of the panels.
  \item \textbf{A persistent error.} With first-order moving-average errors, such as an opening
        error that persists into the next day's bar, the same instruments gave \thn{p6-bad-beta1}
        and \thn{p6-bad-beta2}, and $J$ rejected in a share of \thn{p6-power}.
  \item \textbf{Deeper lags.} Instruments dated $t-2$ restored \thn{p6-fix-beta1} and
        \thn{p6-fix-beta2}.
\end{itemize}
The two-session-lag sensitivity reported in Section~6.6 is therefore the relevant check of
identification, not an optional one.

\subsection{Verification}\label{sec:verification}

Every closed form, identity and inequality above was checked by
\texttt{scripts/\allowbreak 45\_theory\_checks.py} against one of three references:
\begin{itemize}
  \item a Monte Carlo simulation, in which Brownian extremes were drawn exactly from the Brownian
        bridge within each step and the closed form had to lie within four standard errors;
  \item a numerical integral, with a tolerance of $10^{-6}$;
  \item exact arithmetic, for identities and inequalities.
\end{itemize}
The ledger, \texttt{output/tables/table121\_theory\_checks.csv}, has \thn{checks-n} rows, of which
\thn{checks-pass} pass and \thn{checks-fail} fail (Table~\ref{tab:checks}).

\begin{table}[h]
\centering\small
\caption{Checks of the propositions (ledger: \texttt{table121\_theory\_checks.csv}).}\label{tab:checks}
\begin{tabular}{ccc}
\toprule
Proposition & Checks & Passed \\
\midrule
\tabinput{generated/tab_checks.tex}
\bottomrule
\end{tabular}
\end{table}
